\chapter{Dental Procedures Guide}
\index{Dental Procedures Guide}

\section*{Definition and Overview}
Dental procedures are the treatments dentists use to prevent, diagnose, and treat problems of the teeth, gums, and mouth. They range from simple preventive care, such as cleaning and fluoride application, to restorative treatments such as fillings and crowns, and more complex procedures such as root canal treatment, tooth extraction, implants, and orthodontics. Understanding the range of dental procedures helps people know what to expect and make informed decisions about their oral health. This chapter provides a practical guide to common dental procedures and how they are performed.

\section*{Epidemiology}
Dental procedures are among the most common healthcare interventions worldwide. The majority of adults have had a filling, and many have had extractions, crowns, or root canal treatment at some point. Dental implants and orthodontic treatment are increasingly common. In many countries, most dental procedures are performed in private practices, with public services providing subsidised care to eligible groups. Advances in materials, anaesthesia, and technology have made procedures faster, more comfortable, and more durable.

\section*{Aetiology and Pathophysiology}
\begin{itemize}
    \item \textbf{Prevention:} professional cleaning, fluoride treatment, and fissure sealants prevent decay and gum disease
    \item \textbf{Restoration:} fillings, crowns, bridges, and implants restore teeth damaged by decay, trauma, or loss
    \item \textbf{Endodontics:} root canal treatment removes infected pulp from inside the tooth
    \item \textbf{Surgery:} extractions, including wisdom teeth, and surgical placement of implants
    \item \textbf{Orthodontics:} braces and aligners move teeth to correct alignment and bite
    \item \textbf{Periodontics:} scaling, root planing, and gum surgery treat gum disease
    \item \textbf{Prosthodontics:} dentures and other replacements restore missing teeth
\end{itemize}

\section*{Clinical Features}
\begin{itemize}
    \item The type of procedure depends on the condition being treated
    \item Most procedures are performed under local anaesthesia, so they are painless during treatment
    \item Mild discomfort, swelling, or bleeding after procedures usually settles quickly
    \item Some procedures need aftercare, such as soft food, rest, or medication
    \item Follow-up visits may be needed to complete or review treatment
    \item Discussing the procedure, costs, and alternatives with the dentist helps manage expectations
\end{itemize}

\section*{Differential Diagnosis}
\begin{itemize}
    \item Choice between restoring and extracting a tooth
    \item Filling versus crown versus onlay for damaged teeth
    \item Implant versus bridge versus denture for missing teeth
    \item Root canal treatment versus extraction for infected teeth
    \item Orthodontic treatment versus monitoring for mild misalignment
\end{itemize}

\section*{When to Seek Medical Care}
Regular dental check-ups identify the need for procedures early, when treatment is simpler and less costly. See a dentist for any dental problem, and ask about the options, benefits, and costs of recommended treatment. Some procedures require referral to specialists such as oral surgeons, endodontists, or orthodontists.

\textbf{Contact your local emergency services for:}
\begin{itemize}
    \item Severe swelling, difficulty breathing or swallowing, or signs of spreading infection after a procedure
    \item Uncontrolled bleeding that does not stop with pressure
    \item Serious trauma or collapse during or after treatment
\end{itemize}

\section*{Diagnosis and Investigations}
\begin{itemize}
    \item \textbf{Dental examination:} determines which procedures are needed
    \item \textbf{X-rays and scans:} assess decay, bone, roots, and surgical planning
    \item \textbf{Impressions or digital scans:} used for crowns, bridges, and aligners
    \item \textbf{Medical history:} review of health conditions and medicines that affect treatment, such as blood thinners
    \item \textbf{Cost estimate and consent:} the dentist explains the procedure, risks, and costs before treatment
\end{itemize}

\section*{Management}
\begin{itemize}
    \item \textbf{Preventive procedures:} cleaning, fluoride, and sealants
    \item \textbf{Restorations:} fillings, crowns, inlays, and onlays
    \item \textbf{Endodontics:} root canal treatment
    \item \textbf{Extractions:} removal of decayed, infected, or impacted teeth
    \item \textbf{Implants:} surgical placement of artificial roots and crowns
    \item \textbf{Orthodontics:} braces and clear aligners
    \item \textbf{Prosthetics:} dentures, bridges, and implant-supported restorations
    \item \textbf{Periodontal treatment:} scaling, root planing, and gum surgery
    \item \textbf{Aftercare:} following the dentist's instructions, taking prescribed medicines, and attending follow-up
\end{itemize}

\section*{Complications}
\begin{itemize}
    \item Pain, swelling, and bleeding after procedures
    \item Infection after surgery or root canal treatment
    \item Damage to nearby teeth, nerves, or sinuses during complex procedures
    \item Failure of restorations, implants, or other treatments over time
    \item Allergic reactions to materials or medicines
    \item Complications of sedation or anaesthesia, which are rare
    \item Need for further treatment if a procedure fails
\end{itemize}

\section*{Prognosis and Outcomes}
Modern dental procedures have excellent success rates. Fillings, crowns, implants, root canal treatment, and orthodontics all have high long-term success when performed well and followed by good home care. Complications are uncommon, and most can be managed. Choosing the right procedure, with realistic expectations and regular review, gives the best chance of a durable, comfortable, and attractive result.

\section*{Prevention and Self-Care}
\begin{itemize}
    \item Attend regular check-ups so procedures are needed less often
    \item Ask questions about any recommended procedure, its alternatives, and costs
    \item Follow pre-treatment advice, such as avoiding food before sedation
    \item Follow aftercare instructions carefully
    \item Take prescribed medicines and complete antibiotic courses
    \item Maintain good hygiene to protect the results of treatment
    \item Return for scheduled follow-up visits
\end{itemize}

\section*{Sources and Further Reading}
The following reputable sources provide further information for healthcare students and patients seeking reliable, current advice about dental procedures.

\begin{itemize}
    \item Australian Dental Association. \textit{Dental treatments: what to expect.} https://www.ada.org.au/
    \item NHS. \textit{Dental treatments and procedures.} https://www.nhs.uk/conditions/dental-treatments/
    \item Mayo Clinic. \textit{Oral and dental procedures.} https://www.mayoclinic.org/
    \item Healthdirect Australia. \textit{Dental treatments and services.} https://www.healthdirect.gov.au/
    \item MedlinePlus. \textit{Dental procedures.} https://medlineplus.gov/
    \item Oral Health Foundation. \textit{Guide to dental treatments.} https://www.dentalhealth.org/
\end{itemize}
